\scanpage{189}
de diviseurs de $\underline O_{C,c}$.

Posons
\[
\begin{aligned}
S&=\deleted{\operatorname{Spec}}C_a=\operatorname{Spec}\underline O_{C,a}\\
\widehat S&=\operatorname{Spec}\widehat{\underline O_{C,a}}
\end{aligned}
\]
nous avons un hom canonique
\[
\operatorname{Pic}(S-a)\longrightarrow\operatorname{Pic}(\widehat S-\hat a)
\]
[qui, dans le cas $S$ normal, s'interprète comme
\[
\operatorname{Cl}(S)\longrightarrow\operatorname{Cl}(\widehat S)\qquad ]
\]
et on peut se demander si c'est un isomorphisme. \deleted{Faisons écla} Pour ceci, nous allons interpréter l'hom précédent d'une autre manière. Soit $\widetilde C$ obtenu à partir de $C$ en ``faisant éclater l'origine'', donc $\widetilde C$ est un fibré vectoriel sur $X\simeq(\widetilde C)_a$, défini \deleted{\unsure{par des fibrés}} \inserted{comme $W(\underline O_X(-1))$}.

L'on a \deleted{$\unsure{\operatorname{Pic}(C_a-a}$}
\diagram{\begin{tikzcd}[ampersand replacement=\&]
\widetilde C\arrow[d]\&\widetilde C\times_{C_a}C=\widetilde S\arrow[l]\arrow[d]\\
C\&C_a=S\arrow[l]
\end{tikzcd}}
\diagram{\begin{tikzcd}[ampersand replacement=\&,column sep=small,row sep=small]
\operatorname{Pic}(C_a-a)\arrow[d,"\unsure{\sim|2|\wr}"]\&
\operatorname{Pic}(C-a)\arrow[l,"\sim"']\arrow[d,phantom,"\rotatebox{90}{$\simeq$}" description]\&
\deleted{\operatorname{Pic}(\widetilde C-X)}\arrow[l,"\deleted{\sim}"']\arrow[d,phantom,"\deleted{\simeq}" description]\\
\operatorname{Pic}(\widetilde S-X)\&\operatorname{Pic}(\widetilde C-X)\arrow[l,"\sim"']\&
\deleted{\operatorname{Pic}(\widetilde C)/\mathbb Z\,\unsure{p|D|\eta}}
\end{tikzcd}}
% archive sheet number: 188
